\section{Controlled Effects of Exponent Allocation}
\label{sec:controlled-findings}

We begin with interventions that distinguish exponent shape, sampled range,
and the weights learned under an allocation. The Cosh quantiles used here
are constructed in \S\ref{sec:construction}, with endpoint anchoring keeping
their range fixed; an exponential deformation supplies a second family.

\subsection{Controlled effect at fixed range and range interaction}
\label{sec:exp-identify}

We first ask whether interior allocation affects learned behavior when the sampled range is exactly fixed. We train paired $151.9$M models with $K=32$ rotary pairs at sequence length $256$. Within each of three seeds, architecture, initialization, token order, optimizer, schedule, and the $499{,}974{,}144$-token budget are identical. The geometric arm uses the FMRoPE training-base rule $b=L_{\mathrm{train}}$ \citep{oka2026fmrope}; the Cosh arm anchors its $\tau=4$ quantiles to the same endpoints. Only the $30$ interior exponents differ.

At the fixed training support, the mean Cosh-minus-geometric tail-NLL differences at lengths $256, 512, 1024$, and $2048$ are $+0.026, -0.281, -0.176, -0.146$. All three seeds improve at each extended length (Fig.~\ref{fig:spectral-budget-overview}). Interior allocation thus changes learned generalization without any scalar base shift. However, retargeting the log-frequency span to each evaluation length reverses the ordering to $+0.060, +0.227, +0.460$ across all three seeds. This reversal demonstrates that allocation cannot be judged independently of operating range: favorable fixed-support orderings do not survive arbitrary range dilation (Appendix~\ref{sec:identification-details}).

\subsection{Cross-shape factorial across configurations}
\label{sec:factorial}

A $50.9$M factorial varies two bases, two training lengths, three head dimensions, and three paired seeds, retaining endpoints throughout. At the reference Cosh strength, the weighted OOD objective improves in $7/12$ configurations; preassigned $1.25\times$ improves $10/12$, and matched exponential improves $9/12$. The aggregate reference-Cosh-minus-exponential contrast is small and uncertain ($+0.00074$ NLL, $95\%$ interval $[-0.006, 0.008]$; Table~\ref{tab:m4}). These results establish allocation effects across configurations while bounding claims about the unique superiority of any single shape.


\subsection{Weight compatibility and table crossings}
\label{sec:coadaptation-finding}

The frequency table defines a basis in which query and key projections are learned. In a $50$M pair with perplexities near $7$, installing Cosh frequencies into Geo-trained weights raises PPL to $76.20$ and the reverse swap to $23.05$ (Fig.~\ref{fig:weight-table-crossing}). A separate crossing at $151.9$M repeats the effect across two training seeds using derived factor-four runtime tables: each frozen weight set strictly prefers the table derived from its own training allocation. This distinguishes learned compatibility from static basis geometry, without claiming that frequency modifications are universally harmful.

\begin{figure}[t]
\centering
\includegraphics[width=0.92\linewidth]{fig_weight_table_crossing.pdf}
\caption{\textbf{Weights co-adapt to their exponent allocation.}
Each row freezes one trained weight set and changes only its runtime table.
Left: $50$M crossing, PPL on eight windows. Right: $151.9$M crossing, mean
tail NLL on $32$ anchors in each of two training seeds, using extensions
derived from the two training tables. FMR is the geometric table at the fixed
training range; shading shows the NLL increase relative to the matched
diagonal cell.}
\label{fig:weight-table-crossing}
\end{figure}
